
(1.7) que @@M0@@ est
déterminé par sa valeur sur l'élément @@M1@@ de @@M2@@ (resp.\ \ldots)..\ Supposons que @@M3@@ soit un foncteur tel que @@M4@@, etc..~; l'élément @@M5@@ est une série formelle à coefficients dans @@M6@@~; le @@M7@@-ième coefficient de cette série sera appelé \uline{@@M8@@-ième polynôme universel de} @@M9@@~; la connaissance de ces polynômes équivaut donc à la connaissance de @@M10@@. D'après la démonstration de 1.6, ces polynômes sont eux-mêmes déterminés par l'élément @@M11@@, le calcul s'effectuant selon la méthode de la démonstration : on calcule les éléments @@M12@@, en décomposant @@M13@@ en @@M14@@ et en utilisant l'additivité de @@M15@@~; les polynômes ainsi obtenus sont symétriques par rapport aux @@M16@@ (@@M17@@ fixe et @@M18@@ variable) et s'expriment en fonction des @@M19@@~; en recollant grâce à la filtration de @@M20@@, on détermine ainsi les polynômes universels de @@M21@@. Si @@M22@@, le calcul de @@M23@@ se fait en substituant les coefficients @@M24@@ des @@M25@@ aux @@M26@@ dans les polynômes universels.

Supposons qu'il existe un entier @@M27@@ tel que le coefficient de @@M28@@ dans @@M29@@ soit un polynôme dont les composantes isobares ont un poids variant entre @@M30@@ et @@M31@@~; alors le @@M32@@-ième polynôme universel de @@M33@@ a des composantes isobares dont le poids varie entre @@M34@@ et @@M35@@~: ceci résulte immédiatement du mode de calcul exposé plus haut, et du théorème des fonctions symétriques.

